\subsection{2.1+2.2, Vectors, vector equations, and spans}
\begin{problems}
\problem % Dan's notes
True or false: For any collection of vectors $v_1, \ldots, v_k$ in ${\bf R}^n$, the zero vector in ${\bf R}^n$ is a linear combination of $v_1, \ldots, v_k$.

\begin{solution}
True: $0 = 0 v_1 + \cdots + 0 v_k$.
\end{solution}

\problem % Dan's notes
True or false: The vector $\vec{1 1}$ can be written as a linear combination of $\vec{2 2}$ and $\vec{-2 -2}$ in infinitely many ways.

\begin{solution}
 True: the equation $x_1 \vec{2 2} + x_2 \vec{-2 -2} = \vec{1 1}$ corresponds to the augmented matrix $\amat{2 -2 1; 2 -2 1}$.  One step of row-reduction gives us $\amat{2 -2 1; 0 0 0}$, so the system is consistent and has a free variable, so it has infinitely many solutions.  
\end{solution}

\problem % Dan's notes
True or false: The vector $\vec{8 16 1}$ in the span of $\vec{1 2 6}$ and 
$\vec{-1, -2, -1}$.

\begin{solution}
True: row reduction quickly gives $\amat{1 -1 8; 2 -2 16; 6 -1 1} \to \amat{1 -1 8; 0 0 0; 0 5 -47}$.  We don't need to keep going to RREF to see that this augmented matrix has a pivot in every column except the rightmost column, so it is consistent.
\end{solution}

\problem % Dan's notes
True or false: If $v_1, \ldots, v_k$ are vectors in ${\bf R}^n$, then the equation $x_1 v_1 + \cdots + x_k v_k = 0$ must be consistent.

\begin{solution}
True: If we set $x_1$ through $x_k$ all equal to $0$, then 
$0v_1 + \cdots + 0v_k = 0$.  
\end{solution}

\problem % Dan's notes
True or false: The plane $z=1$ in ${\bf R}^3$ is a span.

\begin{solution}
False: any span contains the zero vector but the plane $z=1$ does not contain the zero vector, so it is not a span.
\end{solution}

\problem %Jankowski exam 1 Fall 2017
True or false: The three vectors
  $v_1 = \vec{1 0 0}, \ v_2 = \vec{1 1 0},$ and $v_3 = \vec{-1 0 1}$
  span ${\bf R}^3$. 

\begin{solution}
True.  If $b$ is any vector in ${\bf R}^3$, then the
vector equation $x_1 v_1 + x_2 v_2 + x_3 v_3 = b$ corresponds to the augmented matrix	
\[\amat{1 1 -1 \vert; 0 1 0 b; 0 0 1 \vert} \xrightarrow[]{RREF}
      \amat{1 0 0 *; 0 1 0 *; 0 0 1 *}.\] 
From the RREF we see that the augmented matrix cannot have a pivot in the rightmost column, so the system is consistent no matter what $b$ is.  In other words, every $b$ in ${\bf R}^3$ is in $\Span\{v_1, v_2, v_3\}$.
\end{solution}

\problem %Jankowski exam 1 Fall 2017 
True or false: If $u$, $v$, and $w$ are nonzero vectors in ${\bf R}^2$, 
then Span$\{u, v, w\}$ is ${\bf R}^2$. 

\begin{solution}
False.  Take $v_1=v_2=v_3=\vec{1 0}$.  Many other counterexamples 
possible.    
\end{solution}

\problem % Created for ILA
True or false:  
  The line $y = 2x + 1$ is a span of some number of vectors in
  ${\bf R}^2$.

\begin{solution}
False.  The span of any number of vectors must include the origin, but the 
line $y = 2x + 1$ does not go through the origin.    
\end{solution}

\problem % Dan's notes
Suppose $v$ is a vector in ${\bf R}^n$, and consider the set of all linear combinations of $v$.  What geometric shape(s) can this be?

\begin{solution}
If $v$ is the zero vector, then the set of all linear combinations of $v$ is just the zero vector, since $c(0)=0$ no matter what the scalar $c$ is.

If $v$ is not the zero vector, then the set of all linear combinations of $v$ is a line.  More specifically, it is the line through the origin and $v$.
\end{solution}
% Needs solution picture of vector in 2-d and 3-d?

\problem Consider the following two vectors which characterize two colors in terms of Red, Green, and Blue. % Dan's notes
\[\vec{180 50 200}, \vec{100 150 100}\]
\begin{center}
  \includegraphics[scale=0.6]{2_1-colors-1.png}  
\end{center}
For which $h$ can we get $\vec{116 130 h}$ as a combination of those two colors?
\begin{center}
\includegraphics[scale=0.6]{2_1-colors-2.png}
\end{center}
\begin{solution}
We put everything into an augmented matrix and row reduce.
\begin{align*}
\amat{180 100 116; 50 150 130; 200 100 h} 
&
\xrightarrow[\textup{then } R_1 = \frac{R_1}{50}]{R_1 \leftrightarrow R_2}
\amat{1 3 13/5; 180 100 116; 200 100 h}
\xrightarrow[R_3 = R_3 - 200 R_1]{R_2 = R_2 - 180 R_1}
\amat{1 3 13/5; 0 -440 -352; 0 -500 h-520}
\\
& \xrightarrow[]{R_3 = R_3 - (500/440)R_2}
\amat{1 3 13/5; 0 -440 -352; 0 0 h-120}.
\end{align*}
This system is consistent if and only if $h = 120$.
\end{solution}



\problem % Created for ILA
Consider the system of equations given by
\[5x_1 - x_2 + 5x_3 = 2\]
\[-x_1 - 8x_3 = 4\]
\[2x_1 + x_2 + x_3 = 0.\]
\begin{subproblems}
  \subproblem Write the system of equations as a vector equation.
  \subproblem Solve the vector equation, and use your solution to
  write the target vector in your equation as a linear combination
  of the other vectors.
  %write $b$ as a linear combination of $v_1, v_2, v_3$.
\end{subproblems}

\begin{solution}
  \begin{subproblems}
    \subproblem
    $x_1 \vec{5 -1 2} + x_2 \vec{-1 0 1} + x_3 \vec{5 -8 1}
    = \vec{2 4 0}.$
    \subproblem We row-reduce:
    \[\amat{5 -1 5 2; -1 0 -8 4; 2 1 1 0}
      \xrightarrow[]{RREF} \amat{1 0 0 4/5; 0 1 0 -1; 0 0 1 -3/5}.\]
    This gives $x_1 = 4/5$, $x_2 = -1$, and $x_3 = -3/5$. Therefore,
    \[\frac{4}{5} \vec{5 -1 2} - \vec{-1 0 1}
      - \frac{3}{5} \vec{5 -8 1} = \vec{2 4 0}.\]
  \end{subproblems}
\end{solution}

\problem % Created for ILA
\begin{subproblems}
  \subproblem Write nonzero vectors $v_1$ and $v_2$ in ${\bf R}^2$ so
  that $\textup{Span}\{v_1, v_2\}$ is a line.%
  \subproblem Write nonzero vectors $v_1$ and $v_2$ in ${\bf R}^2$ so
  that $\textup{Span}\{v_1, v_2\}$ is the line $y = -4x$ (where we use
  the coordinates $(x, y)$ to label points in ${\bf R}^2$).
\end{subproblems}

\begin{solution}
  \begin{subproblems}
    \subproblem Many examples possible.  For example
    \[v_1 = \vec{1 1}, \quad v_2 = \vec{2 2}.\] %
    \subproblem We can choose any two nonzero vectors that are on the
    line $y = -4x$, which is the span of $\vec{1 -4}$.  For example,
    \[v_1 = \vec{1 -4}, \quad v_2 = \vec{2 -8}.\]
  \end{subproblems}
\end{solution}

\problem
Let $v_1 = \vec{1 0 0}$, $v_2 = \vec{0 0 1}$, $v_3 = \vec{1 0 1}$.
Describe $\textup{Span}\{v_1, v_2, v_3\}$.

\begin{solution}
  The vector $v_3$ is a linear combination of $v_1$ and $v_2$ since
  $v_1 + v_2 = v_3$, so $\textup{Span}\{v_1, v_2, v_3\}$ is the same
  as $\textup{Span}\{v_1, v_2\}$.  The span of $v_1$ and $v_2$ is the
  $xz$-plane of ${\bf R}^3$, since if $x_1$ and $x_2$ are any real
  numbers then $x_1 v_1 + x_2 v_2 = \vec{x_1 0 x_2}$.  For an
  illustration, see this {\color{blue}
    \href{https://textbooks.math.gatech.edu/ila/1553/demos/spans.html?v1=1,0,0&v2=0,0,1&target=1,0,1&range=3}{graphic}}.
\end{solution}

\problem % Jankowski Spring 2018 exam 1A
Find all solutions to the vector equation
\[x_1 \vec{2 -4 1} + x_2 \vec{4 1 -2} = \vec{2 14 -7}.\] If there are
no solutions, justify why the vector equation is inconsistent.

\begin{solution}
\begin{align*}
  \amat{2 4 2; -4 1 14; 1 -2 -7}
  & \xrightarrow[]{R_1=R_1/2} \amat{1 2 1; -4 1 14; 1 -2 -7} 
\xrightarrow[R_3=R_3-R_1]{R_2=R_2+4R_1} \amat{1 2 1; 0 9 18; 0 -4 -8}
  \\  \\
  & \xrightarrow[R_3=R_3/-4]{R_2=R_2/9} \amat{1 2 1; 0 1 2; 0 1 2} 
 \xrightarrow[R_1=R_1-2R_2]{R_3=R_3-R_2} \amat{1 0 -3; 0 1 2; 0 0 0}.
\end{align*}
Therefore, $x_1= \ -3$ \ \ and \ $x_2=2$.
\end{solution}

\problem % Created for ILA
Supppose we are asked to find a solution to the linear system of
equations corresponding to the augmented matrix
\[\amat{1 -2 3 -1; 2 0 1 3; -1 2 3 4}.\]
\begin{subproblems}
  \subproblem Rewrite the problem in terms of a vector equation.
  \subproblem Rewrite the problem as a question about linear
  combinations.
% \subproblem Rewrite the problem as a question about spans.
\end{subproblems}

\begin{solution}
  \begin{subproblems}
    \subproblem Find a solution to the vector equation
    \[x_1 \vec{1 2 -1} + x_2 \vec{-2 0 2} + x_3 \vec{3 1 3} = \vec{-1
        3 4}.\] %
    \subproblem Write $\vec{-1 2 3 4}$ as a linear combination of
    $\vec{1 2 -1}$, $\vec{-2 0 2}$, and $\vec{3 1 3}$.
  \end{subproblems}
\end{solution}

\problem % Created for ILA
Find all values of $b$ (if there are any) so that $\vec{b -3 2}$ is a
linear combination of $\vec{3 1 2}$ and $\vec{1 0 4}$, and write
$\vec{b -3 2}$ as a linear combination of $\vec{3 1 2}$ and
$\vec{1 0 4}$.

\begin{solution}
  We set up an augmented matrix and solve for $b$ so that the
  corresponding system is consistent.  We will swap the first and
  third rows to make the math easier, but also to make it easier to
  solve for $b$.
  \begin{align*}
    \amat{3 1 b; 1 0 -3; 2 4 2}
    \xrightarrow[\textup{then } R_1 = R_1 / 2]{R_1 \leftrightarrow R_3}
    &
      \amat{1 2 1; 1 0 -3; 3 1 h}
      \xrightarrow[R_3 = R_3 - 3R_1]{R_2 = R_2 - R_1}
      \amat{1 2 1; 0 -2 -4; 0 -5 h-3}
    \\
    \xrightarrow[]{R_2 = (-1/2)R_2}
    &
      \amat{1 2 1; 0 1 2; 0 -5 h-3}
      \xrightarrow[]{R_3 = R_3 + 5R_2}
      \amat{1 2 1; 0 1 2; 0 0 h+7}.
  \end{align*}
  The system is consistent when there is no pivot in the right column,
  so $h+7=0$, thus $h = -7$.  To finish the problem, we need to solve
  \[x_1 \vec{3 1 2} + x_2 \vec{1 0 4} = \vec{-7 -3 2}.\] For this, we
  use the work that was already done above and do one more row
  operation.
  \begin{align*}
    \amat{1 2 1; 0 1 2; 0 0 0}
    \xrightarrow[]{R_1 = R_1 - 2R_2}
    \amat{1 0 -3; 0 1 2; 0 0 0}, \qquad x_1 = -3, \ x_2 = 2.
  \end{align*}
  \[-3 \vec{3 1 2} + 2 \vec{1 0 4} = \vec{-7 -3 2}.\]  
\end{solution}

\problem % Created for ILA
Let $v_1 = \vec{1 0 2}$ and $v_2 = \vec{0 1 3}$.  Write a vector $w$
in ${\bf R}^3$ that is \emph{not} in $\textup{Span}\{v_1, v_2\}$.

\begin{solution}
  Since $\textup{Span}\{v_1, v_2\}$ is the set of all linear
  combinations of $v_1$ and $v_2$, we need to find a vector $w$ that
  is not a linear combination of $v_1$ and $v_2$.  In other words, we
  need to find a $w = \vec{a b c}$ so that the vector equation
  $x_1v_1 + x_2 v_2 = w$ is inconsistent.  We can put this in
  augmented matrix form and row-reduce.
  \[\amat{1 0 a; 0 1 b; 2 3 c}
    \xrightarrow[R_3 = R_3 - 3R_2]{R_3 = R_3 - 2 R_1} \amat{1 0 a; 0 1
      b; 0 0 c-2a-3b}.\] %
  The system will be inconsistent precisely when the final column has
  a pivot, which is when $c - 2a - 3b \neq 0$.  For example, the
  vector $\vec{1 1 1}$ will make the system inconsistent since
  $1-2(1)-3(1) = -4 \neq 0$.
  
  Alternatively, we could have done this problem through trial and
  error.  Since the system of equations is
  $x_1 = a, \ x_2 = b, \ 2x_1 + 3x_2 = c$, one vector that will make
  this inconsistent is $\vec{a b c} = \vec{0 0 1}$, since that
  requires $x_1 = 0$ and $x_2 = 0$ but $2x_1 + 3x_2 = 1$.
\end{solution}

\problem Find all values of $c$ and $d$ so that
$\textup{Span}\left\{\vec{1 c}, \vec{d, 2}\right\}$ is not
${\bf R}^2$.

\begin{solution}
  Neither of the vectors given is the zero vector, so their span
  is at least a line.  The only way their span is just a line is if
  they are collinear.  Therefore, the vectors will span ${\bf R}^2$
  unless one is a scalar multiple of the other, which means there is
  some (nonzero) $k$ so that $d = 1k$ and $2 = ck$, which means
  \[k = d = 2/c.\]
  Therefore, $c$ is any nonzero real number and $d = 2/c$.
\end{solution}

\problem Consider the set $V$ of all vectors in ${\bf R}^3$ of the form
\[\vec{3x_1-x_2 2x_1+5x_2 -10x_2} \quad (x_1, \ x_2
  \textup{ real}).\] Write $V$ as a span of vectors in ${\bf R}^3$.

\begin{solution}
  Unwrapping the general form of a vector in $V$, we see
  \[\vec{3x_1-x_2 2x_1+5x_2 -10x_2} =
    \vec{3x_1 2x_1 0} + \vec{-x_2 5x_2 -10x_2}
    = x_1 \vec{3 2 0} + \vec{-1 5 -10}.\]
  Therefore, the vectors in $V$ are simply all linear combinations
  of $\vec{3 2 0}$ and $\vec{-1 5 -10}$.  In other words,
  \[V = \textup{Span}\left\{ \vec{3 2 0}, \vec{-1 5 -10}.\right\}.\]
\end{solution}

\end{problems}\setcounter{problem}{0}

\newpage

\subsection{2.3, Matrix equations}
\begin{problems} 
\problem % 
True or false: Suppose $A$ is a $2 \times 4$ matrix.  Then the equation
  $Ax = b$ must be inconsistent for some $b$ in ${\bf R}^2$.
 
\begin{solution}
False.  Since $A$ has $2$ rows and $4$ columns, it is possible for $A$ to 
have a pivot in each row, in which case $Ax=b$ will be consistent for 
every $b$ in ${\bf R}^2$.  For example:    
\[A = \mat{1 0 0 0; 0 1 0 0}.\]
\end{solution}

\problem % 
True or false: Suppose $A$ is an $m \times n$ matrix and the matrix 
equation $Ax = b$ is consistent for each $b$ in ${\bf R}^m$. 
Then $m \leq n$.
 
\begin{solution}
True.  If $m > n$, then $A$ has more rows than columns
    and so the max number of pivots possible number of pivots is $n$,
    thus $A$ cannot have a pivot in every row. 
\end{solution}

\problem % Dan's notes
True or false: If $A$ is a $3 \times 5$ matrix and the product $Ax$ makes sense, then which ${\bf R}^n$ is $x$ in?

\begin{solution}
The vector $x$ must be in ${\bf R}^5$, since $A$ has $5$ columns.
\end{solution}

\problem % Dan's notes
 Write the following linear equation as a matrix equation:
  \[x + y + z = 1\]
  \[x - 2y - z = 5.\]

\begin{solution}
$\mat{1 1 1; 1 -2 -1} \vec{x y z} = \vec{1 5}$.
\end{solution}

\problem
Write the following linear system as a matrix equation:
    \begin{align*}
        x_1   + 2x_2 + 7x_4                      &= 10 \\
        3x_1  + 7x_2 + \tfrac{1}{2}x_3           &= 7  \\
        x_2  + x_3 + 15x_4                       &= 30 \\
        5x_1  +  12x_3                           &= 17
    \end{align*}

\begin{solution}
    \(
    \mat{1 2 0 7; 3 7 \tfrac{1}{2} 0; 0 1 1 15; 5 0 12 0}
    \vec{x_1 x_2 x_3 x_4}
    =
    \vec{10 7 30 17}
    \)
\end{solution}

\problem  
Write the following linear system as a matrix equation:  
\[
\begin{aligned}
2x_1 - x_2 + 3x_3 &= 5 \\
x_1 + 4x_2        &= -2 \\
3x_2 + 2x_3       &= 7
\end{aligned}
\]

\begin{solution}
\(
\mat{2 -1 3; 1 4 0; 0 3 2}
\vec{x_1 x_2 x_3}
=
\vec{5 -2 7}
\)
\end{solution}

\problem  
Write the following linear system as a matrix equation:  
\[
\begin{aligned}
x_1 + 3x_2 &= 12 \\
-2x_1      &= -6 \\
4x_2       &= 8 \\
5x_1 + x_2 &= 3
\end{aligned}
\]

\begin{solution}
\(
\mat{1 3; -2 0; 0 4; 5 1}
\vec{x_1 x_2}
=
\vec{12 -6 8 3}
\)
\end{solution}

\problem  
Write the following linear system as a matrix equation:  
\[
\begin{aligned}
x_1 + 2x_3 - x_4 + 6x_5 &= 14 \\
3x_2 - 5x_3 + x_4        &= -2
\end{aligned}
\]

\begin{solution}
\(
\mat{1 0 2 -1 6; 0 3 -5 1 0}
\vec{x_1 x_2 x_3 x_4 x_5}
=
\vec{14 -2}
\)
\end{solution}





\problem % Dan's notes
Carry out the following matrix multiplication.
\[\mat{0 2; 0 4; 5 0} \vec{3 2}.\]
\begin{solution}
\[\mat{0 2; 0 4; 5 0} \vec{3 2} 
= 3 \vec{0 0 5} + 2 \vec{2 4 0} = \vec{0 0 15} + \vec{4 8 0}
= \vec{4 8 15}.\]  
\end{solution}

\problem % Created for ILA
Carry out the following matrix multiplication.
\[\mat{2 3 -5; -3 1 0} \vec{3 -1 2}.\]
\begin{solution}
  \begin{align*}
    \mat{2 3 -5; -3 1 0} \vec{3 -1 2}
    & = 3 \vec{2 -3} + (-1) \vec{3 1} + 2 \vec{-5 0} \\
    & = \vec{6 -9} + \vec{-3 -1} + \vec{-10 0} \\
    & = \vec{-7 -10}.
  \end{align*}
\end{solution}

\problem % Dan's notes
Is the following matrix equation consistent?  If so, solve it.
\[\mat{1 1; 1 0; 1 1} \vec{x_1 x_2} = \vec{2 3 2}.\]
Can you do this problem without row-reducing?

\begin{solution}
This matrix equation is consistent: $x_1 = 3$ and $x_2 = -1$.  We can do this without row-reduction.  We see $x_1$ must equal $3$ in order for the second entry to be $3$, and $3\vec{1 1 1} = \vec{3 3 3}$ so we need to subtract $\vec{1 0 1}$ to get $\vec{2 3 2}$.
\end{solution}

\problem % Dan's notes
Is the following matrix equation consistent?  If so, solve it.
\[\mat{1 1; 1 0; 1 1} \vec{x_1 x_2} = \vec{2 3 3}.\]
Can you do this problem without row-reducing?

\begin{solution}
 This matrix equation is inconsistent.  Each of the vectors $\vec{1 1 1}$ and $\vec{1 0 1}$ has its first entry equal to its third entry (so any linear combination will also have first and third entries equal), but the target vector $\vec{2 3 3}$ has first entry different than third entry.   
\end{solution}


  
\problem % Created for ILA
Todd Stromboly sells apples, bananas, and dragonfruits.  At his fruit
stand:

You purchased $2$ apples, $4$ bananas, and $1$ dragonfruit, at a cost
of $\$3$.

Your friend bought $3$ apples, $2$ bananas, and $2$ dragonfruits at a
cost of $\$4$.

Your sworn enemy bought $1$ apple, $2$ bananas, and $5$ dragonfruits
at a cost of $\$6$.  \newline Let $a$, $b$, and $c$, be the sale price
of an apple, banana, and dragonfruit, respectively.
\begin{subproblems}
  \subproblem Write a system of equations whose solution would give us
  the price of each item.  %
  \subproblem Write a vector equation whose solution would give us the
  price of each item.
  %
  \subproblem Write a matrix equation whose solution would give us the
  price of each item.  %
  \subproblem Find the price of each item.
\end{subproblems}
	%
\begin{solution}
  \begin{subproblems}
	\subproblem The system of equations is
	\[2a + 4b + c = 3\]
	\[3a + 2b + 2c = 4\]
	\[a + 2b + 5c = 6.\]
    %
    \subproblem The vector equation is
	\[a \vec{2 3 1} + b \vec{4 2 2} + c \vec{1 2 5} = \vec{3 4 6}.\]
    %
	\subproblem The matrix equation is
	\[\mat{2 4 1; 3 2 2; 1 2 5}x = \vec{3 4 6}.\]
    %
	\subproblem We put an augmented matrix into RREF.
	\[\amat{2 4 1 3; 3 2 2 4; 1 2 5 6} \xrightarrow[]{RREF}
      \amat{1 0 0 1/2; 0 1 0 1/4; 0 0 1 1}.\] Thus $a = \frac{1}{2}$,
    \ $b = \frac{1}{4}$, and $c=1$. In other words, an apple costs
    $\$0.50$, a banana costs $\$0.25$, and a dragonfruit costs $\$1$.
	\end{subproblems}	
\end{solution}
	
\problem Suppose $A = \mat{| | | |; v_1 v_2 v_3 v_4; | | | |}$ is a
matrix whose columns are nonzero vectors and
  \[5v_1 - v_2 + v_3 - v_4 = 0.\]
  \begin{subproblems}
	\subproblem Write one nonzero vector $b$ so that $Ax=b$ is
    consistent.  \subproblem Write one nonzero solution to the
    equation $Ax=0$.
	\end{subproblems}
	%
    \begin{solution}
      \begin{subproblems}
        \subproblem Many possible answers.  For example, $Ax = v_1$ is
        consistent since
        \[A \vec{1 0 0 0} = \mat{| | | |; v_1 v_2 v_3 v_4; | | | |}
          \vec{1 0 0 0} = v_1.\] %
        \subproblem By definition of matrix multiplication,
        \[\mat{| | | |; v_1 v_2 v_3 v_4; | | | |}
          \vec{5 -1 1 -1} = 5v_1 - v_2 + v_3 - v_4 = 0\] so
        $x= \vec{5 -1 1 -1}$ satisfies $Ax=0$.
	\end{subproblems}
  \end{solution}


\problem % Created for ILA
%
Your overall semester score in Math $9123$ will be determined by
the weights: $15\%$ homework, $15\%$ quizzes, $20\%$ for midterm 1,
$20\%$ for midterm 2, and $30\%$ for the final exam.
\newline \newline
Your grades were $95\%$ on homework, $85\%$ for quizzes, $83\%$
for midterm 1, and $87\%$ for midterm 2.  
\begin{subproblems}
  \subproblem Write $A$, $x$, and $b$ so that solving the matrix
  equation $Ax = b$ will give us the final exam score we need in order
  to get a $90\%$ overall score for the course.  %
  \subproblem Solve the equation in part (a).
\end{subproblems}

\begin{solution}
  \begin{subproblems}
    \subproblem We arrange the matrix entries of $A$ in the order that
    the categories were given, and we put the grades in each category as
    the corresponding entries of $x$.
    \[A = \mat{0.15 0.15 0.20 0.20 0.30}, \qquad
      x = \vec{95 85 83 87 x_5}, \qquad b = 90.\]
    \subproblem Carrying out the multiplication and solving gives us
    \[(0.15)95 + (0.15)85 + (0.20)83 + (0.20)87 + (0.30)x_5 = 90\]
    \[61 + 0.30 x_5 = 90, \qquad 0.30 x_5 = 29, \qquad x_5 =
      290 / 3 = 96.\bar{6}\]
    If the exam does not give fractions of a point as credit, you will
    need to get a $97\%$ on the final exam.
  \end{subproblems}
\end{solution}

\problem % Jankowski exam 1 Fall 2017
    The diagram below represents the temperature at points along
    wires, in celcius.
\begin{center}
  \begin{tikzpicture}[scale=1, thick, every node/.style={inner
      sep=3pt, label distance=1mm}]
    %\node at (0,2.4) {Traffic flow (cars/hr)};
   \node at (0,2.3) {$30$};
    \node at (1.9,-2) {$60$};
    \node at (-1.9,-2) {$90$};
    \node at (-0.4,1) {$T_1$};
    \node at (1.4,-1) {$T_2$};
     \node at (-1.4,-1) {$T_3$};
    \point[scale=1.5] (A) at (0,1);
    \point[scale=1.5] (C) at (1,-1);
    \point[scale=1.5] (D) at (-1,-1);
     \point[scale=1.5] (P) at (0,2);
     \point[scale=1.5] (Q) at (3/2,-2);
     \point[scale=1.5] (R) at (-1.5,-2);
   % \draw[mid arrow] (A.center) to["$x_1$"] (C.center);
   \draw(A.center) to (C.center);
    %\draw[mid arrow] (D.center) to["$x_2$"] (C.center);
    \draw(D.center) to (C.center);
    %\draw[mid arrow] (D.center) to["$x_3$"] (A.center);
    \draw(D.center) to (A.center);
    %\draw[rmid arrow=.3] (A.center) to["$110$"] +(-1/2,1);
    \draw(A.center) to+(0,1);
     \draw(C.center) to+(1/2,-1);
    \draw(D.center) to+(-1/2,-1);
  \end{tikzpicture}

\end{center}

Let $T_1, \ T_2, \ T_3$ be the temperatures at the interior points.
Assume the temperature at each interior point is the average of the
temperatures of the three adjacent points.
\begin{subproblems}
  \subproblem Write a system of three linear equations whose solution
  would give the temperatures $T_1$, $T_2$, and $T_3$.  Do not solve
  it.  %
  \subproblem Write the system as a vector equation. Do not
  solve it.  %
  \subproblem Write a matrix equation $Ax=b$ that represents this
  system.  Specify every entry of $A$, $x$, and $b$.  Do not solve it.
\end{subproblems}

\begin{solution}
\begin{subproblems}
  \subproblem The left side system below or right-side system below
  are both fine.

  $T_1 = \dfrac{T_2+T_3+30}{3}$, \quad or \quad
  $3T_1 - T_2 - T_3 = 30$. \newline $T_2 = \dfrac{T_1+T_3+60}{3}$,
  \quad or \quad $-T_1+3T_2-T_3 = 60$. \newline
  $T_3 = \dfrac{T_1+T_2+90}{3}$, \quad or \quad $-T_1-T_2+3T_3 = 90$.
  \newline

  \subproblem
  $T_1 \vec{3 -1 -1} + T_2 \vec{-1 3 -1} + T_3 \vec{-1 -1 3} = \vec{30
    60 90}$.  \newline

  \subproblem
  $\mat{3 -1 -1; -1 3 -1; -1 -1 3} \vec{T_1 T_2 T_3} = \vec{30 60
    90}$.
\end{subproblems}
\end{solution}

 \end{problems}\setcounter{problem}{0}
\newpage

\subsection{2.4, Solution sets}
\begin{problems}
\problem % Created for ILA
True or false: If $A$ is a $4 \times 6$ matrix with four pivots, then
  the matrix equation $Ax = b$ must be consistent for every $b$ in
  ${\bf R}^4$.  %
  
\begin{solution}
True.  Since $A$ has $4$ pivots, it follows that for
    any vector $b$ in ${\bf R}^4$, the matrix $(A \ | \ b)$ will have
    a pivot in each of its $4$ rows, and each pivot will be to the
    left of the augment bar.  Therefore, the system $Ax = b$ will be
    consistent.  %
\end{solution}

\problem % Created for ILA
If $A$ is a $3 \times 5$ matrix with three pivots, then
  the set of solutions to $Ax = 0$ is a plane. %
  
\begin{solution}
True.  Since $A$ has $3$ pivots but $5$ columns, there
    will be exactly $2$ columns that do not have a pivot.  Therefore,
    the system $Ax = 0$ will have exactly two free variables, which
    means that its solution set is a plane.
\end{solution}

\problem % Jankowski Spring 2018 exam 1
True or false: If $b$ is a nonzero vector and $Ax=b$ has infinitely many 
solutions, then $Ax=0$ must have infinitely many solutions.

\begin{solution}
True.  The solution set to $Ax=b$ is a translate of the solution set to 
$Ax=0$, hence $Ax=0$ must also have infinitely many solutions. 
\end{solution}

\problem  % Jankowski Spring 2018 exam 1
True or false: If $A$ is a $3 \times 4$ matrix and the solution set for 
$Ax=0$ is a line, then $A$ has exactly $2$ pivots.  \newline %
\begin{solution}
False.  The solution set to $Ax=0$ will be a line if and only if $A$ has 
exactly $1$ free variable.  Since $A$ has $4$ columns, this means $A$ has 
exactly $3$ (not $2$) pivots. %
\end{solution}

\problem %
True or false: There is a $3 \times 2$ matrix $A$ whose homogeneous solution set
$Ax=0$ is the line $x_1 + x_2 = 1$. %

\begin{solution}
False.  The zero vector $x_1 = x_2 = 0$ does not satisfy the equation 
$x_1 + x_2 = 1$.
\end{solution}


\problem %Dan's notes
Suppose that the set of solutions to $Ax = b$ is the plane $z=1$.  What is the
set of solutions to $Ax = 0$?

\begin{solution}
  From the theory of section 2.4, the solution set to $Ax = 0$ is the plane
  through the origin parallel to the plane $z=1$.  Therefore, the solution set
  to $Ax = 0$ is the $x_1x_2$-plane in ${\bf R}^3$ (this is the plane $z=0$).
\end{solution}

\problem %Dan's notes
Suppose that the set of solutions to $Ax = 0$ is the plane
$y=x$ in ${\bf R}^2$.  Is it possible that there is a $b$ so that
the set of solutions to $Ax = b$ is the line $x+y=1$?  Justify your answer.

\begin{solution}
No.  If $Ax=b$ is consistent, then its solution set must be parallel to the solution set to $Ax=0$, namely the line $y=x$ which has slope $1$.  The line $x+y=1$ is the line $y= -x + 1$, which has slope $-1$, so it cannot be the solution set to any consistent system $Ax=b$.
\end{solution}

\problem %Dan's notes
Suppose that the set of solutions to $Ax = 0$ is the plane
$x+y = 0$ in ${\bf R}^3$.  Is it possible that there is a $b$ so that
the set of solutions to $Ax = b$ is the $z$-axis in ${\bf R}^3$?  Justify your answer.

\begin{solution}
  No.  The solution set to $Ax=0$ is a plane, then the solution set of any
  consistent equation $Ax=b$ must be a parallel plane, but the $z$-axis is a
  line.
\end{solution}

\problem %Dan's notes
Suppose that the set of solutions to $Ax = 0$ is the plane
$x + 2y - 3z = 0$ in ${\bf R}^3$ and that the vector $\vec{1 3 5}$ is a solution to $Ax = b$. 
\begin{subproblems}
    \item Find one other solution to $Ax=b$.
    \item Find all solutions to $Ax=b$.  Write your answer in parametric vector form.
\end{subproblems}

\begin{solution}
\begin{subproblems}
\subproblem We can get the solution set to $Ax=b$ by taking any particular solution to $Ax = b$ (in this case, $\vec{1 3 5}$) and then adding all homogeneous solutions.  The solution set to $Ax=0$ is the plane $x+2y-3z=0$, and we can find a homogeneous solution by just guessing and checking to find $x$, $y$, and $z$ that satisfy the equation.  One such homogeneous solution is $\vec{3 0 1}$.  Another solution to $Ax = b$ is
\[\vec{1 3 5} + \vec{3 0 1} = \vec{4 3 6}.\]
\subproblem The homogeneous equation $x+2y-3z=0$ corresponds to $\amat{1 2 -3 0}$, which gives $x = -2y + 3z$ where $y$ and $z$ are free.
\[\vec{x y z} = \vec{-2y+3z y z} = \vec{-2y y 0} + \vec{3z 0 z}
= y \vec{-2 1 0} + z \vec{3 0 1}.\]
The solution set to $Ax = b$ is the translation of the homogeneous solution set by the particular solution $\vec{1 3 5}$:
\[\vec{1 3 5} + y \vec{-2 1 0} + z \vec{3 0 1}, \quad \textup{ where } y, z \textup{ real}.\]
\end{subproblems}    
\end{solution}

\problem %Dan's notes
Is there a $2 \times 2$ matrix so that the set of solutions to $Av = \vec{1 2}$ is the line $y = x + 1$?  If so, find such an $A$.  If not, explain why not.

\begin{solution}
Yes: we need $\vec{1 2}$ to be in the span of the columns of $A$, so that $Av = \vec{1 2}$ is consistent.  Also, we need for the line $y=x$ to be the set of solutions to $Ax=0$ since that is the line through the origin parallel to $y=x+1$.  One way to get $A$ is to start as if both columns are $\vec{1 2}$, then modify $A$ until we satisfy the above.  One such matrix is \[A = \mat{-1 1; -2 2}.\]
\end{solution}

\problem % Created for ILA
Fill in the blanks: If $A$ is a $5 \times 3$ matrix with two pivots,
then the set of solutions to $Ax = 0$ is a \underline{\qquad \quad}
that lives in \underline{\qquad}.

\begin{solution}
  $A$ has two pivot columns, so we will get one free variable in the
  when we solve $Ax = 0$, therefore the solution set will be a line.
  The solution set lives in ${\bf R}^3$ since $A$ has three columns.
\end{solution}


\problem % Created for ILA
Let
\[A = \mat{1 -2 1; 0 1 2}, \qquad
  b_1 = \vec{1 -1 0}, \qquad
  b_2 = \vec{2 3}.\]
\begin{subproblems}
  \subproblem Which matrix equation makes sense: $Ax = b_1$ or
  $Ax = b_2$? %
  \subproblem For the answer you chose in part (a), write the matrix
  equation as a system of equations. %
  \subproblem For the answer you chose in part (a), write the matrix
  equation as a vector equation. %
  \subproblem Solve the matrix equation and write your answer in
  parametric vector form.
\end{subproblems}

\begin{solution}
  \begin{subproblems}
    \subproblem $A$ is a $2 \times 3$ matrix, so any solution $x$ must
    be in ${\bf R}^3$ and the target vector must be in ${\bf R}^2$.
    Therefore, $Ax = b_2$ makes sense but $Ax = b_1$ does not. One way
    to check this is to translate the matrix equation into a system of
    equations to confirm our answer, which we do in part (b). %
    \subproblem The linear system corresponding to $Ax = b_2$ is
    \[x_1 - 2x_2 + x_3 = 2\]
    \[0x_1 + x_2 + 2x_3 = 3.\] %
    \subproblem The corresponding vector equation is
    \[x_1 \vec{1 0} + x_2 \vec{-2 1} + x_3 \vec{1 2} = \vec{2 3}.\] %
    \subproblem We row-reduce $(A \ | \ b_2)$:
    \[\amat{1 -2 1 2; 0 1 2 3} \xrightarrow[]{R_1 = R_1 + 2R_2}
      \amat{1 0 5 8; 0 1 2 3}.\] This gives us $x_1 + 5 x_3 = 8$ and
    $x_2 + 2x_3 = 3$, with $x_3$ free.
    \[x_1 = 8 - 5x_3, \qquad x_2 = 3 - 2x_3, \qquad x_3 = x_3
      \textup{ (free)}.\]
    We finish the problem by putting this into parametric vector form:
    \[\vec{x_1 x_2 x_3} = \vec{8-5x_3 3-2x_3 x_3}
      \ = \ \vec{8 3 0} + \vec{-5x_3 -2x_3 x_3}
      \ = \ \vec{8 3 0} + x_3 \vec{-5 -2 1}.\]
  \end{subproblems}
\end{solution}

  
\problem %Jankowski Spring 2018 exam 1
  Consider the system of equations in $x_1$, $x_2$, $x_3$, and $x_4$
  given below.
    \[x_1 - x_2 + 2x_3 - 2x_4 \ = \ -1\]
    \[-x_1 + x_2 -2x_3 + x_4 = \ 2\]
    \[-4x_1 + 4x_2 -8x_3 + 6x_4 \ = \ 6.\] Put an augmented matrix
    into reduced row echelon form to solve the system of equations.
    Write your answer in parametric vector form.

    \begin{solution}
      \begin{align*}
        \amat{1 -1 2 -2 -1; -1 1 -2 1 2; -4 4 -8 6 6}
        &
          \xrightarrow[R_3=R_3+4R_1]{R_2=R_2+R_1}
          \amat{1 -1 2 -2 -1; 0 0 0 -1 1; 0 0 0 -2 2}
          \xrightarrow[R_3= R_3/-2]{R_2= \ -R_2}
          \amat{1 -1 2 -2 -1; 0 0 0 1 -1; 0 0 0 1 -1}
        \\ \\
        & \xrightarrow[R_1= R_1+2R_2]{R_3=R_3-R_2}
          \amat{1 -1 2 0 -3; 0 0 0 1 -1; 0 0 0 0 0}
      \end{align*}

      The free variables are $x_2$ and $x_3$:
      \[x_1= x_2-2x_3-3, \qquad x_2=x_2, \qquad x_3=x_3, \qquad x_4=
        -1.\]

      \[\vec{x_1 x_2 x_3 x_4} = \vec{x_2-2x_3-3 x_2 x_3 -1}
        = \vec{x_2 x_2 0 0} + \vec{-2x_3 0 x_3 0} + \vec{-3 0 0 -1} \
        = \ \boxed{\vec{-3 0 0 -1} + x_2 \vec{1 1 0 0} + x_3 \vec{-2
          0 1 0}}.\]
\end{solution}

\problem Consider the linear system of equations in $x_1$, $x_2$,
$x_3$, and $x_4$ given below.
\[-2x_1 - 2x_2 + x_3 - 4x_4 = -8\]
\[-x_1 - x_2 + x_3 - 2x_4 = -5\]
\[4x_1 + 4x_2 - x_3 + 8x_4 = 14.\]
\begin{subproblems}
  \subproblem Write the set of solutions to the system in parametric
  vector form.
  %
  \subproblem Write two different solutions to the system of equations.
  %
  \subproblem Write two \emph{nonzero} solutions to the corresponding
  \emph{homogeneous} linear system of equations.
\end{subproblems}
\begin{solution}
\begin{subproblems}
  \subproblem We put the system into an augmented matrix and
  row-reduce:
  \[\amat{-2 -2 1 -4 -8; -1 -1 1 -2 -5; 4 4 -1 8 14}
    \xrightarrow[]{RREF}
    \amat{1 1 0 2 3; 0 0 1 0 -2; 0 0 0 0 0}.\]
  Therefore, $x_1 = 3 - x_2 - 2x_4$ and $x_3 = -2$ while $x_2$ and
  $x_4$ are free:
  \[\vec{x_1 x_2 x_3 x_4} = \vec{3-x_2-2x_4 x_2 -2 x_4}
    = \boxed{\vec{3 0 -2 0} + x_2 \vec{-1 1 0 0} + x_4 \vec{-2 0 0 1}}.
  \]
  %
  \subproblem We can obtain a solution to the system by choosing any
  values we wish for $x_2$ and $x_4$.  For example, if we choose $x_2 = 0$
  and $x_4 = 0$ we get
  \[\vec{x_1 x_2 x_3 x_4} = \vec{3 0 -2 0}.\]
  If we choose $x_2 = 1$ and $x_4 = 0$ we get
  \[\vec{x_1 x_2 x_3 x_4} = \vec{2 1 -2 0}.\]
  %
  \subproblem The solutions to the corresponding homogeneous system
  (where $-8$, $-5$, and $14$ are all replaced by zero) are given by
  everything in the span of the two vectors $\vec{-1 1 0 0}$ and
  $\vec{-2 0 0 1}$.  Therefore, $\vec{-1 1 0 0}$ and $\vec{-2 0 0 1}$
  are two solutions to the homogeneous equations.
\end{subproblems}  
\end{solution}


\problem % Created for ILA
Suppose $A$ is a $2 \times 2$ matrix and that for some vector $b$, the
set of solutions to $Ax = b$ is the line $x_2 = 4-x_1$.  Write the
general solution to the homogeneous equation $Ax=0$ in parametric
vector form.

\begin{solution}
  The solution set for $Ax = 0$ must include the origin and also must
  be a translate to the solution set to $Ax = b$.  In other words,
  the solution set for $Ax = 0$ is the line through the origin that
  is parallel to $x_2 = 4-x_1$.  This is the line $x_2 = -x_1$,
  so $x_1 = -x_2$ and $x_2$ is free.
  \[\vec{x_1 x_2} = \vec{-x_2 x_2} = x_2 \vec{-1 1}.\]
  Another way to do this problem is to draw the line $x_2 = 4-x_1$,
  then draw its parallel line through the origin and notice that this
  line is the span of $\vec{-1 1}$ (or equivalently the span of
  $\vec{1 -1}$).
\end{solution}

\problem % Created for ILA, inspired by Gerandy Brito
For some matrix $A$ and some vector $b$, the dashed line below
is the solution set to the matrix equation $Ax = b$.
\begin{center}
  \includegraphics[scale=0.6]{./2_4/solution_sets1.pdf}
\end{center}
\begin{subproblems}
  \subproblem On the same graph, draw the solution set to $Ax = 0$.
  %
  \subproblem True or false: It is possible that for some vector $v$,
  the solution set to $Ax = v$ is a single point.
\end{subproblems}

\begin{solution}
  \begin{subproblems}
    \subproblem The solution set is the line through the origin that
    is parallel to the dashed blue line.
    \begin{center}
      \includegraphics[scale=0.6]{./2_4/solution_sets1-sol.pdf}
    \end{center}
    %
    \subproblem False.  If $Ax = v$ is consistent, then its solution
    set must be a translation of the solution set to $Ax = 0$, which
    is a line.  If $Ax = v$ is inconsistent, then it has no solution.
    Either way, it is not possible for $Ax = v$ to have exactly one
    solution.
  \end{subproblems}
\end{solution}


%\problem % Created for ILA
%Suppose that the set of solutions to some linear system of equations
%in the variables $(x_1, x_2, x_3, x_4)$ has the following parametric
%vector form:
%\[\vec{x_1 x_2 x_3 x_4}
%  = \vec{2 -1 1 0} + x_2 \vec{-3 1 0 0} + x_4 \vec{5 0 0 1}.\]

%\problem Suppose a 


\problem % Created for ILA
Write a system of linear equations whose solution set has parametric
form
\[x = \vec{3 0} + x_2 \vec{2 1}.\]

\begin{solution}
  Writing $x$ as a single vector, we get:
  \[x = \vec{3 0} + \vec{2x_2 x_2} = \vec{3+2x_2 x_2}.\] Therefore,
  $x_1 = 3+x_2$ and $x_2 = x_2$ (free).  One possible answer is the
  following single equation, which will give us $x_1 = 3+x_2$ with $x_2$
  as a free variable:
  \[x_1 - x_2 = 3.\]
\end{solution}

\problem  % Fall 2017 sample exam
The diagram below describes traffic in a part of town.
\begin{center}
  \begin{tikzpicture}[scale=1.4, thick,
      every node/.style={inner sep=3pt, label distance=1mm}]
    \node at (0,2.4) {Traffic flow (cars/hr)};
    \point[scale=1.5] (A) at (0,1);
    %\point[scale=1.5] (B) at (1,1);
    \point[scale=1.5] (C) at (1,-1);
    \point[scale=1.5] (D) at (-1,-1);
  %  \draw[mid arrow] (A.center) to["$x_1$"] (B.center);
    \draw[mid arrow] (A.center) to["$x_1$"] (C.center);
    \draw[mid arrow] (D.center) to["$x_2$"] (C.center);
    \draw[mid arrow] (D.center) to["$x_3$"] (A.center);
    \draw[rmid arrow=.3] (A.center) to["$110$"] +(-1/2,1);
    \draw[mid arrow=.5] (A.center) to["$40$" swap] +(1/2,1);
   % \draw[mid arrow=.7] (B.center) to["$70$" swap] +(1,0);
   % \draw[rmid arrow=.3] (B.center) to["$120$" swap] +(0,1);
   % \draw[mid arrow=.7] (C.center) to["$80$"] +(1,0);
    \draw[mid arrow=.7] (C.center) to["$320$"] +(1,0);
    \draw[mid arrow=.7] (C.center) to["$40$"] +(0,-1);
    \draw[rmid arrow=.7] (D.center) to["$200$" swap] +(-1,0);
    \draw[rmid arrow=.3] (D.center) to["$90$" swap] +(0,-1);
  \end{tikzpicture}

\end{center}
\begin{subproblems}
\item Write a system of three linear equations in $x_1$, $x_2$, and
  $x_3$ corresponding to the traffic flow.
\item Use an augmented matrix to solve this system of linear equations.
\item Were we given enough information to know the exact 
  values of $x_1$, $x_2$, and $x_3$?  If your answer is yes, justify why.
  If your answer is no, give the range of possible values for $x_1$,
  $x_2$, and $x_3$.
\end{subproblems}


\begin{solution}
 \begin{subproblems}  
   \subproblem For the top, bottom right, and bottom left nodes, the
   number of cars entering must match the number of cars exiting, so
   the system is:
   \[x_1 + 40 = x_3 + 110\]
   \[x_1 + x_2 = 360\]
   \[x_2 + x_3 = 290.\] %
   \subproblem The system can be written
   \[x_1 - x_3 = 70\]
   \[x_1 + x_2 = 360\]
   \[x_2 + x_3 = 290.\] We form an augmented matrix and perform row
   operations.
   \[\amat{1 0 -1 70; 1 1 0 360; 0 1 1 290}
     \xrightarrow[]{R_2 = R_2 - R_1} \amat{1 0 -1 70; 0 1 1 290; 0 1 1
       290} \xrightarrow[]{R_3 = R_3 - R_2} \amat{1 0 -1 70; 0 1 1
       290; 0 0 0 0}.\] Therefore, $x_3$ is a free variable,
   $x_1 = x_3 + 70$, and $x_2 = 290-x_3$. \newline \newline

   \subproblem We cannot know the exact values of $x_1$, $x_2$, and
   $x_3$ with the information we have only been given.  For example,
   we could have $x_3=0$, $x_2 = 290$, $x_1=70$.  Or, we could have
   $x_3= 100$, $x_2= 190$, $x_1=170$, etc. \newline \newline
 %
   Even though $x_3$ is a free variable in the system of equations,
   the real-world constraints of the problem make us consider what is
   actually possible here.  Since the variables are counting the
   number of cars at each intersection, they must all be at least
   zero.  Therefore, we need to have $x_3 \geq 0$ (whereby
   $x_1 = x_3 + 70$), and since $x_2 = 290-x_3 \geq 0$ we also need
   $x_3 \leq 290$.%
   \newline \newline
 %
   To summarize: As the values of $x_3$ range from $0$ to $290$, the
   values of $x_1$ range from $70$ to $360$ while the values for $x_3$
   range from $290$ (when $x_3=0$) down to $0$ (when $x_3 = 290$).
 \end{subproblems}
\end{solution}


%\problem Suppose $A$ is a 


\end{problems} \setcounter{problem}{0}

\newpage

\subsection{2.5, Linear independence}

\begin{problems}
\problem % perhaps ILA but perhaps from somewhere else
True or false: Suppose $v_1$ and $v_2$ are vectors in ${\bf R}^2$ 
and neither vector is a scalar multiple of the other.  Then 
$\{v_1, v_2\}$ must be a linearly independent set. %

\begin{solution}
True.  The vectors $v_1$ and $v_2$ cannot be
      linearly dependent, since then $x_1 v_1 + x_2 v_2 = 0$ would
      have a solution where $x_1 \neq 0$ or $x_2 \neq 0$.  This would
      mean $v_1$ and $v_2$ are scalar multiples of each other: if
      $x_1 \neq 0$ then $v_1 = - \frac{x_2}{x_1} v_2$, and similarly
      if $x_2 \neq 0$ then $v_2 = -\frac{x_1}{x_2} v_1$.  %
      %
\end{solution}

\problem % perhaps ILA but perhaps from somewhere else
True or false: %
Suppose $v_1$, $v_2$, and $v_3$ are vectors in ${\bf R}^3$ and 
that $v_j$ and $v_k$ are not collinear whenever $j \neq k$. 
$\Span\{v_j, v_k\}$ is a plane whenever $j \neq k$.  Then 
$\{v_1, v_2, v_3\}$ is a linearly independent set.  %

\begin{solution}
False.  Let $v_1 = \vec{1 0 0}$, \ $v_2 = \vec{0 1 0}$, \ 
$v_3 = \vec{1 1 0}$.  No two among them are collinear, so any two vectors 
chosen among them will span a plane.  However, they are linearly dependent 
since $v_1 + v_2 = v_3$.  %
\end{solution}

\problem % perhaps ILA but perhaps from somewhere else
True or false: %
If $A$ is $4 \times 2$ matrix and the solution set to $Ax=\vec{1 2 2 1}$ 
is a line, then the columns of $A$ are linearly dependent. %

\begin{solution}
True.  From section 2.4, we know that the solution set of any consistent 
system $Ax=b$ is a translation of the homogeneous solution set.  Here, the 
system $Ax=\vec{1 2 2 1}$ has a line as its solution set, so the
homogeneous system $Ax=0$ has a line as its solution set and
thus has infinitely many solutions.  Therefore, the columns of $A$ are 
linearly dependent.   
\end{solution}

\problem % perhaps ILA but perhaps from somewhere else
True or false: %
If a set $S$ of vectors contains fewer vectors than there are entries in 
the vectors, then the set is linearly independent. %

\begin{solution}
False.  For example, the set $S=\left\{\vec{1 0 0}, \vec{0 0 0}\right\}$ 
in ${\bf R}^3$ is linearly dependent.   
\end{solution}

\problem % Dan's notes
True or false: %
If $u$ and $v$ are vectors in ${\bf R}^5$, then the set 
$\{u, v, \sqrt{2}u - \pi v\}$ is linearly independent.

\begin{solution}
False.  The third vector is a linear combination of the first two vectors, 
so the set is linearly dependent.
\end{solution}

\problem % Dan's notes
True or false: %
If a matrix $A$ has a pivot in each column, then the columns of $A$ are 
linearly independent.

\begin{solution}
True.  Having a pivot in each column guarantees that $Ax = 0$ has only the 
trivial solution, which is what it means for the columns of $A$ to be 
linearly independent.
\end{solution}

\problem %Dan's notes
If you have a set of linearly independent set $S$ of vectors, and the span of
the vectors in $S$ is a line, then how many vectors are in $S$?

\begin{solution}
  There is exactly one vector in $S$.  The span of any single nonzero vector is
  a line, whereas a linearly independent set with two (or more) vectors would
  span at least a plane.
\end{solution}
    
\problem % Dan's notes
In each case answer Always, Sometimes, or Never.  Justify your answers.
\begin{subproblems}
  \subproblem A collection of 99 vectors in ${\bf R}^{100}$ is linearly
  dependent. %
  \subproblem A collection of 100 vectors in ${\bf R}^{99}$ is
  linearly dependent.
\end{subproblems}

\begin{solution}
  \begin{subproblems}
    \subproblem False.  We could get a set of 99 linearly independent vectors
    $\{v_1, v_2, \ldots, v_{99}\}$ by putting $1$ in the $k^{th}$ entry of $v_k$
    for each $k$ and making all other entries $0$.
    \[v_1 = \vec{1 0 0 \vdots \, 0 0}, \quad v_2 = \vec{0 1 0 \vdots \, 0 0},
      \cdots, \quad v_{99} = \vec{0 0 0 \vdots \, 1 0}.\] \subproblem True.  If
    we put the 100 vectors as columns of a matrix, we will get a $99 \times 100$
    matrix.  This means the matrix has $100$ columns but a max of $99$ pivots,
    so it cannot have a pivot in each column, therefore the set of vectors is
    linearly dependent.
    \end{subproblems}
  \end{solution}

    \problem Find all values of $h$ so that the following set vectors is linearly independent.
    \[\left\{ \vec{5 7 1}, \vec{-5 7 0}, \vec{10 0 h}\right\}\]
    
    \problem % Fall 2017 Jankowski
    In each case, determine whether the set of vectors is linearly
    dependent or linearly independent.
    \begin{subproblems}
      \subproblem
      $\left\{\vec{1 0 0}, \ \vec{2, -5, 0}, \ \vec{3, 1, 7}\right\}$
      \subproblem
      $\left\{\vec{1 2 3}, \ \vec{0, 0, -1}, \ \vec{2, 4, 6}\right\}$
      \subproblem
      $\left\{\vec{1 0}, \ \vec{0 1}, \ \vec{-2 3}\right\}$
    \end{subproblems}

    \begin{solution}
      \begin{subproblems}
        \subproblem Forming a matrix $A$ with the vectors as columns
        gives 3 pivot columns (no row-reduction required), so the
        vectors are linearly independent. Alternatively, use the
        Increasing Span Criterion. %
        \subproblem The first and third vectors are multiples of each
        other, so the vectors are linearly dependent.  %
        \subproblem Any three vectors in ${\bf R}^2$ are automatically
        linearly dependent.
      \end{subproblems}
    \end{solution}
 
    \problem Suppose that $v_1, v_2, v_3, v_4, b$ are vectors in
    ${\bf R}^6$ and the equation
    \[x_1 v_1 + x_2 v_2 + x_3 v_3 + x_4 v_4 = b\]
    has exactly one solution.
    \begin{subproblems}
      \subproblem Must $\{v_1, v_2, v_3, v_4\}$ be
      linearly independent?
      %
      \subproblem Must the equation
      \[x_1 v_1 + x_2 v_2 + x_3 v_3 = 0\] have exactly one solution?
      %
      \subproblem Must there be a vector $v$ so that
      $\{b, v\}$ is linearly independent and the equation
      \[x_1 v_1 + x_2 v_2 + x_3 v_3 + x_4 v_4 = v\] is consistent?
    \end{subproblems}
    \begin{solution}
      \begin{subproblems}
        \subproblem Yes.  Since the first equation has exactly one
        solution, the matrix $A$ whose columns are $v_1$ through $v_4$
        has a pivot in every column, therefore
        $\{v_1, v_2, v_3, v_4\}$ are linearly independent.
        %
        \subproblem Yes.  Since $\{v_1, v_2, v_3, v_4\}$ is linearly
        independent, so is $\{v_1, v_2, v_3\}$, which is the same
        thing as saying $x_1 v_1 + x_2 v_2 + x_3 v_3 = 0$ has only
        the trivial solution $x_1 = x_2 = x_3 = 0$.
        %
        \subproblem Yes. Since $\{v_1, v_2, v_3, v_4\}$ is linearly
        independent, there is at least one vector $v$ in
        $\{v_1, v_2, v_3, v_4\}$ that is not a scalar multiple of $b$
        (otherwise they would all be scalar multiples of $b$ and
        therefore linearly dependent).  Since $v$ is in the set
        $\{v_1, v_2, v_3, v_4\}$ it is also in
        $\textup{Span}\{v_1, v_2, v_3, v_4\}$, therefore there is at
        least one solution to 
        \[x_1 v_1 + x_2 v_2 + x_3 v_3 + x_4 v_4 = v.\]
        It will become much easier to see why this part of the problem
        is true by the end of chapter 2.
      \end{subproblems}
    \end{solution}

\end{problems} \setcounter{problem}{0}

\newpage

\subsection{2.6, Subspaces}
\begin{problems}
%a
%
\problem %Dan's notes
True or false: The set of solutions to a consistent matrix equation is 
always a subspace.

\begin{solution}
False.  The set of solutions to a consistent matrix equation $Ax = b$ is 
only a subspace when $b=0$.  If $b$ is not the zero vector, then the 
solution set to $Ax = b$ will not contain the zero vector.
\end{solution}

\problem %Dan's notes
True or false: The set of ${\bf R}^n$ that consists only of the zero 
vector is a subspace of ${\bf R}^n$. 

\begin{solution}
True.  The set $\{0\}$ contains the zero vector, and it is closed under 
addition and scalar multiplication because any linear combination of the 
$0$ vector is itself the zero vector.
\end{solution}

\problem 
True or false: If $A$ is a $9 \times 15$ matrix, then the null space of 
$A$ is a subspace of ${\bf R}^9$.  

\begin{solution}
    False: $\Nul(A)$ is a subspace of ${\bf R}^{15}$.
\end{solution}

\problem 
True or false: Every subspace of ${\bf R}^n$ is the span of some
  vectors in ${\bf R}^n$.

\begin{solution}
True.  It is a fundamental fact that every subspace of
    ${\bf R}^n$ is the span of some vectors in ${\bf R}^n$, and,
    conversely, the span of any number of vectors in ${\bf R}^n$ is a
    subspace of ${\bf R}^n$.
\end{solution}

\problem %Dan's notes
Consider the set $S = \{(x, y) \textup{ in } {\bf R}^2 \ | \ xy \geq 0\}$.  Is
$S$ a subspace of ${\bf R}^2$?  If not, which properties does it fail?

\begin{solution}
  (1) $S$ contains the zero vector, since $0\cdot 0 = 0 \geq 0$.  \newline %
  (2) $S$ is not closed under addition.  For example, take $u = \vec{1 2}$ and
  $v = \vec{-2 -1}$.  Both vectors are in $S$ since $1\cdot 2 > 0$ and
  $(-2)(-1) = 2 >0$.  However, their sum is $u+v = \vec{-1 1}$, which is not in
  $S$ since $(-1)(1)<0$.  \newline %
  (3) $S$ is closed under scalar multiplication.  If $u = \vec{x y}$ is in $S$
  and $c$ is any scalar, then for $cu = \vec{cx cy}$ we have
    \[(cx)(cy) = c^2 (xy) \geq 0 \quad \textup{ since } c^2 \geq 0 \textup{ and } xy \geq 0.\]

    Most students find it easier if we first draw a picture of $S$ to gain
    intuition, rather than plugging in random vectors to test properties.  In
    order for a vector to be in $S$, either both entries are positive, or both
    entries are negative, or at least one entry is zero.  Putting this together,
    we see that $S$ is the union of the first and third quadrants of the
    $xy$-plane (including the $x$-axis and $y$-axis).
\end{solution}

\problem %Dan's notes
Consider the $x$-axis in ${\bf R}^3$.  Is it a subspace of ${\bf R}^3$?  If not,
which properties does it fail?

\begin{solution}
  Yes: the $x$-axis in ${\bf R}^3$ is the span of $\vec{1 0 0}$, and spans are
  automatically subspaces of the appropriate ${\bf R}^n$how, so the $x$-axis is a
  subspace of ${\bf R}^3$.
\end{solution}

\problem %Dan's notes
Consider the set
$S = \{(x, y, z, w) \textup{ in } {\bf R}^4 \ | \ x + y - z + w = 0\}$.  Is $S$
a subspace of ${\bf R}^4$?  If not, which properties does it fail?

\begin{solution}
  Yes: $S$ is the solution set to the augmented system $\amat{1 1 -1 1 0}$,
  which is precisely $\Nul\mat{1 1 -1 1}$.  Null spaces are always subspaces of
  the appropriate ${\bf R}^{n}$, and in this case $S$ is the null space of a
  matrix with $4$ columns, so $S$ is a subspace of ${\bf R}^{4}$.
\end{solution}

\problem Let $W$ be the green graph below.  Determine which properties
of a subspace of ${\bf R}^2$ are satisfied by $W$.
\begin{center}
  \includegraphics[scale=0.7]{./2_6/2_6-lineorigin}
\end{center}
\begin{solution}
  (1) $W$ contains the zero vector, since the line passes through the origin.
  \newline \newline %
  (2) $W$ is closed under addition.  If $v$ and $w$ are vectors along the line
  $W$, their sum is just another vector along the line.  \newline \newline %
  (3) $W$ is closed under scalar multiplication.  If you scale any vector along
  $W$, the result is still a vector along $W$.  \newline \newline %
\end{solution}

\problem Let
$V = \left\{ \vec{a b c} \ \textup{ in } {\bf R}^3 \ : \ a-b = c\right\}$.  Is
$V$ a subspace of ${\bf R}^3$?  Justify your answer.

\begin{solution}
  Yes: $V$ is the set of solutions to the homogeneous system of equations
  $a-b-c = 0$.  In other words, $V = \Nul\mat{1 -1 -1}$.  The null space
  of any matrix with $3$ columns is a subspace of ${\bf R}^3$.

  Alternatively, we could verify the subspace properties individually.
  For example, $V$ contains $\vec{0 0 0}$ since $0 - 0 = 0$, and we
  could similarly verify that $V$ is closed under addition and
  multiplication.
\end{solution}


\problem % Created for ILA, inspired by Joe from Fall 2017
Which of the following are subspaces of ${\bf R}^4$?  Answer yes or nox
in each case.
\begin{subproblems}
  \subproblem $\Col\mat{4 3 2 1; 1 2 3 4; 4 3 2 1}$ %
  \subproblem $\Nul\mat{1 2 3 4; 1 2 3 4; 1 2 3 4}$ %
  \subproblem The set of solutions to the system of equations in the
  variables $x_1, x_2, x_3, x_4$ given by
  \[x_1 - x_2 + x_3 + x_4 = 1\]
  \[x_1 + x_4 = 2.\]
\end{subproblems}
\begin{solution}
\begin{subproblems}
  \subproblem No.  This is a subspace of ${\bf R}^3$, not ${\bf R}^4$.
  %
  \subproblem Yes.  The null space of any $3 \times 4$ matrix is a
  subspace of ${\bf R}^4$. %
  \subproblem No.  This solution set does not contain the origin,
  since it corresponds to a non-homogeneous linear system of
  equations.
\end{subproblems}
\end{solution}

\problem If $V$ and $W$ are sets in ${\bf R}^n$, then their \textit{union},
denoted $V \cup W$, is the set of all vectors that are in at least one of $V$
and $W$.  In other words, a vector $x$ is in $V \cup W$ if $x$ is in $V$ or $x$
is in $W$ (or $x$ is in both $V$ and $W$).

Suppose $U$ and $W$ are subspaces of ${\bf R}^n$.  When will $U \cup W$ also a
subspace of ${\bf R}^n$?

\begin{solution}
  The union will be a subspace of ${\bf R}^n$ if and only if one of one of $U$
  or $V$ is contained in the other (or both are equal).  In this case, the union
  $U \cup V$ will just be one of $U$ or $V$, which we already know to be
  subspaces.

On the other hand, suppose neither $U$ nor $V$ is contained in the other.  Then there is a vector $u$ that is in $U$ but not in $V$, and there is a vector $v$ that is in $V$ but not in $U$.  We show that:
\[u+v \textup{ is not in } U \cup V.\] After all, $u+v$ cannot be in the
subspace $U$, because if it were then $(u+v)-u = v$ would need to be in $U$
since $U$ is closed under addition.  Similarly, $u+v$ cannot be in $V$ because
then $(u+v)-v = u$ would need to be in $V$ since $V$ is closed under addition.

In summary: $u$ and $v$ are in $U \cup V$, but $u+v$ is not in $U \cup V$, so
$U \cup V$ is not closed under addition and is therefore not a subspace of
${\bf R}^n$.
\end{solution}


\problem Let $W$ be the green graph below.  Determine which properties
of a subspace of ${\bf R}^2$ are satisfied by $W$.
\begin{center}
  \includegraphics[scale=0.7]{./2_6/2_6-which.pdf}
\end{center}
\begin{solution}
  $W$ doesn't satisfy any of the properties of a subspace: \newline
  (1) $W$ does not pass through the origin. \newline \newline %
  (2) $W$ is not closed under addition since for example $(-2,0)$ and
  $(0, 2)$ are in $W$ but their sum $(-2,2)$ is not in $W$. \newline
  \newline %
  (3) $W$ is not closed under multiplication since, for example,
  $(-2, 0)$ is in $W$ but $-(-2,0)=(2, 0)$ is not in $W$. \newline
  \newline %
  Although it fails all subspace properties, the truth is that $W$ is
  \emph{almost} a subspace because it is a line in ${\bf R}^2$, but
  the fact that it does not pass through the origin ruins all three
  subspace properties.
\end{solution}

\problem Let $W$ be the red graph below.  It is the paraboloid
$z = x^2 + y^2$. Determine which properties of a subspace of
${\bf R}^3$ are satisfied by $W$.
\begin{center}
  \includegraphics[scale=0.3]{./2_6/paraboloid}
\end{center}
\begin{solution}
  (1) $W$ contains the zero vector since it passes through the
  origin. \newline \newline %
  (2) $W$ is not closed under scalar multiplication because if $u$ is
  any nonzero vector in $W$ then it is above the $z$-axis, but $-u$ is
  below the $z$-axis and thus not in $W$.  \newline \newline %
  (3) $W$ is not closed under addition because, for example, the point
  $(0, 1, 1)$ is in $W$ but $(0, 2, 2)$ is not in $W$.  %
  %In reality, $W$ is the paraboloid $z = x^2 + y^2$.
\end{solution}

\problem Let $W$ be the light blue graph below.  Determine which
properties of a subspace of ${\bf R}^3$ are satisfied by $W$.
\begin{center}
  \includegraphics[scale=0.4]{./2_6/planeno0}
\end{center}
\begin{solution}
  $W$ is a plane in ${\bf R}^3$ that does \emph{not} contain the
  origin.  It will fail all three properties of a subspace in the same
  way that a line not through the origin in ${\bf R}^2$ will fail all
  three subspace properties.
\end{solution}
\problem Recall from section 2.1 that a QR code is a
$29 \times 29$ grid of boxes, where each box is colored black or
white.  We can think of each box as labeled by $1$ (if white) or $0$
(if black).  Reading from top to bottom and left to right, we can
regard a QR code as a vector $x$ in
${\bf R}^{29 \times 29} = {\bf R}^{841}$, where each entry of $x$ is
$0$ or $1$.

Let $S$ be the set of all QR codes, considered as a subset of
${\bf R}^{841}$.  Determine which properties of a subspace $S$
satisfies.
\begin{subproblems}
  \subproblem Does $S$ contain the $0$ vector?  \subproblem Is $S$
  closed under addition?  \subproblem Is $S$ closed under scalar
  multiplication?
\end{subproblems}
	
\begin{solution}
  \begin{subproblems}
    \subproblem Yes.  The vector $v = (0, 0, \ldots, 0)$ is the QR
    code where every box is black. %
      %
    \subproblem No.  The vectors $x = (1, 0, 0, \ldots 0)$ and
    $y=(1, 1, 0, \ldots, 0)$ are in $S$.  However, if we add $x$ and
    $y$ we obtain $x+y = (2, 1, 0, \ldots, 0)$, which is not in $S$
    because it has an entry (its first entry) which is neither $0$ nor
    $1$.
      %
    \subproblem No.  The vector $x = (1, 0, 0, \ldots, 0)$ is in $S$
    but $2x = (2, 0, 0, \ldots, 0)$ is not in $S$ because it has an
    entry (its first entry) which is neither $0$ nor $1$.
  \end{subproblems}
\end{solution}

\problem Let $V$ be the set of all points $(x, y)$ in ${\bf R}^2$
satisfying $x - y > 1$.
\begin{subproblems}
  \subproblem Does $V$ contain the zero vector?
  %
  \subproblem Is $V$ closed under addition?%
  \subproblem Is $V$ closed under scalar multiplication?
\end{subproblems}
\begin{solution}
  \begin{subproblems}
    \subproblem No, since $0-0$ is not greater than $1$.  %
    \subproblem Yes: take any vectors $u = \vec{x_1 y_1}$ and
    $v = \vec{x_2 y_2}$ in $V$, so $x_1 - y_1 \geq 1$.  For the vector
    $u+v = \vec{x_1+x_2 y_1+y_2}$, we subtract the second coordinate
    from the first and regroup terms:
    \begin{align*}
      (x_1+x_2)-(y_1+y_2)
      &  = x_1+x_2 - y_1 - y_2 \\
      &  = (x_1 - y_1) + (x_2 - y_2) \\
      & > 1 + 1 > 1,
    \end{align*}
    therefore $u+v$ is in $V$.
      %
    \subproblem No: for example, $u = \vec{1 0}$ is in $V$, but
    $-u = \vec{-1 0}$ is not in $V$ since $-1 - 0 = -1$ which is not
    greater than $1$.
  \end{subproblems}
  An alternative approach is to draw $V$, which is relatively
  straightforward since $V$ consists of all points satisfying
  $y < x - 1$, which is all points that are strictly below the line
  $y = x-1$.  Once you graph $V$, you can tell right away that it does
  not contain the zero vector and is not closed under scalar
  multiplication, and you can also deduce that it is closed under
  addition.
  \end{solution}

  \problem Let $V$ be the set of all points $(x, y, z)$ in ${\bf R}^3$
  satisfying $xyz \geq 0$.
  \begin{subproblems}
    \subproblem Does $V$ contain the zero vector?
    \subproblem Is $V$ closed under addition?
    \subproblem Is $V$ closed under scalar multiplication?
  \end{subproblems}
  %
  \begin{solution}
    \begin{subproblems}
      \subproblem Yes: the origin satisfies
      $0\cdot 0 \cdot 0 = 0 \geq 0$.%
      \subproblem No: for example, let
      \[u = \vec{1 -1 -1} \qquad v = \vec{1 0 2}.\] Then $u$ and $v$
      are both in $V$ but their sum is $u+v = \vec{2 -1 1}$ which is
      not in $V$ since $(2)(-1)(1) = -2 < 0$.
    \end{subproblems}
  \end{solution}
  
  \problem (Challenge problem) Abstract spaces go beyond just
  ${\bf R}^n$.  For example, the set of all polynomials with real
  coefficients forms an abstract vector space.  \newline \newline %
  Consider the set $W$ consisting of all polynomials that have degree
  $2$ or less.  In other words, $W$ is the set of all functions of the
  form
  \[f(x) = a_0 + a_1 x + a_2 x^2\]
  where $a_0$, $a_1$, and $a_2$ are real numbers.
  Show that $W$ satisfies the properties of a subspace:
  \begin{subproblems}
    \subproblem $W$ contains the zero polynomial.  In other words, $W$
    contains the polynomial given by $p(x) = 0$ for all $x$.
    %
    \subproblem If polynomials $f$ and $g$ are in $W$, then so is the
    polynomial $f+g$ given by \[(f+g)(x) = f(x) + g(x).\]
    %
    \subproblem If $f$ is in $W$ and $c$ is a scalar, then the
    polynomial $cf$ given by $(cf)(x) = cf(x)$ is also in $W$.
  \end{subproblems}

\end{problems}\setcounter{problem}{0}

\newpage

\subsection{2.7, Basis and dimension}

\begin{problems}
  \problem %Dan's notes
  True or false: If $S$ is a plane containing the origin in ${\bf R}^3$, then
  every basis for $S$ has exactly two vectors.

\begin{solution}
  True. Since $S$ is a plane through the origin in ${\bf R}^3$, it is a subspace
  of ${\bf R}^3$ and it is two-dimensional, so any basis for $S$ conssists of
  exactly two vectors.
\end{solution}

\problem %Dan's notes
True or false: If two vectors are in the same plane through the origin in
${\bf R}^3$, and if the two vectors are not collinear, then they form a basis
for the plane.

\begin{solution}
True.  The span of the two vectors must be a plane, otherwise their span 
would be a line or a point which is impossible since the two vectors is 
not collinear.  Any two vectors that span a plane must be linearly 
independent, therefore the two vectors form a basis for the plane.
\end{solution}

\problem %Dan's notes
True or false: The dimension of the null space of $A$ is the number of 
pivots of $A$.

\begin{solution}
False: $\dim(\Nul \ A)$ is the number of columns in $A$ that do NOT have 
a pivot.
\end{solution}

\problem %Dan's notes
True or false: If $b$ lies in the column space of $A$ and the columns of $A$ are
linearly independent, then $Ax = b$ has infinitely many solutions.

\begin{solution}
  False: the columns of $A$ are linearly independent, so $A$ has a pivot in
  every column, therefore the augmented matrix $(A | b)$ has a pivot in each
  column to the left of the augment bar, so either the system is inconsistent or
  has exactly one solution.  We were told that $b$ is in the column space of
  $A$, so the system is consistent, hence $Ax = b$ has exactly one solution.
\end{solution}

\problem % Dan's notes
True or false: Any three vectors that span ${\bf R}^3$ must be linearly 
independent.

\begin{solution}
  True.  Let $A$ be the $3 \times 3$ matrix whose columns are the three vectors.
  Since these vectors span ${\bf R}^3$, we conclude $A$ has a pivot in every
  row, so $A$ has $3$ pivots.  This means that $A$ has a pivot in each of its
  $3$ columns, so its columns are linearly independent.
\end{solution}

\problem How many vectors are there in a basis for a line $S$ through the origin
in ${\bf R}^7$? %Dan's notes

\begin{solution}
  One.  Take any nonzero vector $v$ along the line $S$.  Then $\{v\}$ is
  linearly independent, and any vector along the line $S$ can be obtained by
  scaling $v$ by an appropriate scalar $c$.  In fact, the set of all vectors of
  the form $cv$ (where $c$ is a scalar) is the line $S$. Therefore, $\{v\}$ is a
  basis for $S$.
\end{solution}

\problem Find a basis for the $yz$-plane in ${\bf R}^3$. %Dan's notes

\begin{solution}
  The $yz$-plane is all vectors of the form $\vec{0 y z}$ where $y$ and $z$ are
  real.  A basis for the $yz$-plane is
  $\left\{\vec{0 1 0}, \vec{0 0 1}\right\}$: the two vectors are clearly
  linearly independent, and they span the $yz$-plane since any vector of the
  form $\vec{0 y z}$ can be written as the linear combination
\[\vec{0 y z} = y \vec{0 1 0} + z \vec{0 0 1}.\]
\end{solution}


\problem Find a basis for the column space and the null space of
\[A = \mat{1 2 3; 4 5 6; -7 -8 -9}.\] %Dan's notes

\begin{solution}
We row-reduce $\amat{A 0}$ to find
\[\amat{1 2 3 0; 4 5 6 0; -7 -8 -9 0}
\xrightarrow[]{RREF} 
\amat{1 0 -1 0; 0 1 2 0; 0 0 0 0}.\]
The pivot columns of $A$ form a basis for $\Col(A)$:
\[\left\{\vec{1 4 -7}, \vec{2 5 -8}\right\}.\]
For $\Nul(A)$, the RREF of $\amat{A 0}$ gives us $x_1 - x_3 = 0$, $x_2 + 2x_3 = 0$, and $x_3$ is free, so $x_1 = x_3$ and $x_2 = -2x_3$.
\[\vec{x_1 x_2 x_3} = \vec{x_3 -2x_3 x_3} = x_3 \vec{1 -2 1}.
\qquad \textup{Basis for } \Nul(A): \left\{\vec{1 -2 1}\right\}\]
\end{solution}

\problem Find a basis for ${\bf R}^3$ where no basis vector has $0$ as entry. %Dan's notes

\begin{solution}
Many answers possible, for example:
\[\left\{\vec{1 1 1}, \vec{1 2 1}, \vec{1 1 2}\right\}.\]
We could row-reduce to find that $\mat{1 1 1; 1 2 1; 1 1 2}$ has a pivot in
every row and column. \newline This answer was obtained by starting with the
standard basis
\[\left\{\vec{1 0 0}, \vec{0 1 0}, \vec{0 0 1}\right\}\]
and doing row replacements until there were no zeros left in the matrix (added the second and third rows to the first row, then added the new first row to the second and third rows).
\end{solution}

\problem 
Allen Argent has put the matrix $A$ below in its reduced row echelon
form:
\[A=\mat{1 -2 0 4; -7 14 3 2; 4 -8 -2 -4} \xrightarrow[]{RREF} \mat{1
    -2 0 4; 0 0 1 10; 0 0 0 0}.\]
\begin{subproblems}
  \subproblem Find a basis $\mathcal{B}$ for $\Nul(A)$.
    %
  \subproblem Is $x=\vec{2 3 -10 1}$ in $\Nul(A)$?  If so, write $x$
  as a linear combination of the basis vectors of $\Nul(A)$.  If not,
  justify why $x$ is not in $\Nul(A)$.
    %
  \subproblem Is $\vec{0 -3 2}$ in $\Col(A)$?
\end{subproblems}

\begin{solution}
  \begin{subproblems}
	\subproblem The RREF of $A$ shows that if $Ax=0$ then
	\[x_1 = 2x_2 - 4x_4, \ x_2 = x_2, \ x_3 = -10x_4, \ x_4 = x_4.\]
	So $\vec{x_1 x_2 x_3 x_4} = \vec{2x_2-4x_4 x_2 -10x_4 x_4}
	= x_2 \vec{2 1 0 0} + x_4 \vec{-4 0 -10 1}$.  
	Thus $\mathcal{B} = \left\{\vec{2 1 0 0}, \vec{-4 0 -10 1}\right\}.$
    %
    \subproblem
	\[\amat{2 -4 2; 1 0 3; 0 -10 -10; 0 1 1}
      \xrightarrow[R_3= -R_3/10]{R_1 \leftrightarrow R_2} \amat{1 0 3;
        2 -4 2; 0 1 1; 0 1 1} \xrightarrow[R_4=R_4-R_3]{R_2=R_2-2R_1}
      \amat{1 0 3; 0 -4 -4; 0 1 1; 0 0 0}
      \xrightarrow[R_3=R_3-R_2]{R_2=R_2/-4} \amat{1 0 3; 0 1 1; 0 0 0;
        0 0 0}.\]%
    Therefore, $x$ is in $\Nul(A)$, in fact $x=3b_1 + b_2$.
    %
    \subproblem Yes.  $\vec{0 -3 2}$ is
    $(-1)*(\textup{third column of $A$})$ so it is in $\Col(A)$, no
    work required.
  \end{subproblems}
\end{solution}
%
%
\problem Let $A=\mat{1 -2; -2 4}$.  Clearly draw $\Col(A)$ and
$\Nul(A)$, using the usual convention of labeling points in
${\bf R}^2$ by $(x,y)$.%

\includegraphics[scale=0.6]{./2_7/blank-graph} \qquad
\quad \includegraphics[scale=0.6]{./2_7/blank-graph}
  
 \begin{solution}  
   $\Col(A)$ is the span of $\vec{1 -2}$; $\Nul(A)$ is the span of $\vec{2 1}$.
   We graph $\Col(A)$ first, then graph $\Nul(A)$.  For emphasis, we draw the
   column space a solid line and the null space as a dashed line.
 
 \includegraphics[scale=0.6]{./2_7/2_7-pr-11-col.pdf} \qquad \qquad
 \includegraphics[scale=0.6]{./2_7/2_7-pr-11-null.pdf} 
 \end{solution}

 \problem % Created for ILA
 Write a matrix whose column space is the solid line and whose null space is the
 dashed line graphed below.
\begin{center}
  \includegraphics[scale=0.6]{./2_7/2_7-pr-12.pdf}
\end{center}
\begin{solution}
  Here $\Col(A)$ is the span of $\vec{1 3}$, so each column must be a scalar
  multiple of $\vec{1 3}$.  Also, we need $\Nul(A)$ to be the span of
  $\vec{1 -4}$, so $x_1 = -x_2/4$, which is the homogeneous equation
  $x_1 + x_2/4 = 0$.  Therefore, each row of $\amat{A 0}$ must be a scalar
  multiple of $\amat{1 1/4 0}$.  Some possible examples for $A$ is
\[\mat{1 1/4; 3 3/4}, \qquad \mat{4 1; 12 3}, \qquad
\mat{-4 -1; -12 -3}.\]
\end{solution}
 
\problem Suppose $W$ is a $4$-dimensional subspace of ${\bf R}^{13}$.
Must it be true that $W = \Col(A)$ for some matrix $A$?
Justify your answer.

\begin{solution}
  Yes.  Just from the fact that $\dim(W) = 4$, we know that there is a
  basis of $W$ consisting of exactly $4$ vectors
  $v_1, \ v_2, \ v_3, \ v_4$ so that
  $W = \Span\{v_1, v_2, v_3, v_4\}$.  Therefore, $W = \Col(A)$ for the
  matrix $A = \mat{ | | | | ; v_1 v_2 v_3 v_4; | | | | }.$
\end{solution}
 
\problem %(made for ILA)
Let
 \[V = \left\{\vec{x y z} \textup{ in } {\bf R}^3
     \ : \ x = y \right\}.\]
 \begin{subproblems}
   \subproblem Show that $V$ is a subspace of ${\bf R}^3$.%
   \subproblem Find a basis for $V$.  (Hint: it may be helpful to
   write $V$ as the null space of some matrix) %
   \subproblem Write a matrix $B$ so that $\Col(B) = V$.
 \end{subproblems}


 \begin{solution}
   \begin{subproblems}
     \subproblem We could verify that $V$ is a subspace using the
     properties of a subspace, but there is a much faster way.  $V$ is
     the set of all points $(x, y, z)$ in ${\bf R}^3$ so that
     $x - y = 0$.  This is $\textup{Nul}\mat{1 -1 0}$, which is
     automatically a subspace of ${\bf R}^3$ since the null space of
     \textit{any} matrix with $3$ columns is a subspace of
     ${\bf R}^3$.
     %
     \subproblem To find $\textup{Nul}\mat{1 -1 0}$ we use the
     augmented matrix $\amat{1 -1 0 0}$, which gives us $x_1 = x_2$,
     $x_2$ is free, and $x_3$ is free.  Therefore, a basis for $V$ is
     \[\vec{1 1 0}, \vec{0 0 1}.\]
     %
     \subproblem This problem might seem complicated at first, but if
     we go back to the fundamentals, we see it is a concept check
     using the definition of column space.  We just need to write a
     matrix $B$ so that
     $\Col(B)= \textup{Span}\left\{\vec{1 1 0}, \vec{0 0 1}\right\} =
     V$.  One way to do this is just to let
     \[B = \mat{1 0; 1 0; 0 1}.\]     
   \end{subproblems}
 \end{solution}

 \problem  (Challenge) Consider the set $W$ consisting of all polynomials that have degree
  $2$ or less.  In other words, $W$ is the set of all functions of the
  form
  \[f(x) = a_0 + a_1 x + a_2 x^2\]
  where $a_0$, $a_1$, and $a_2$ are real numbers.
  \begin{subproblems}
    \subproblem What should it mean for a set of polynomials in $W$ to be
    linearly independent? %
    \subproblem Find a basis for $W$.  In other words, find a linearly
    independent set $S$ of polynomials in $W$, so that any polynomial in $W$ can
    be written as a linear combination of elements in $S$.
    %
    \subproblem What is the dimension of $W$?
  \end{subproblems}
  %
  \begin{solution}
    We use ${\bf 0}$ to denote the zero polynomial.  In order words, it is the
    polynomial \[{\bf 0}(x) = 0 + 0 x + 0x^2.\]
  \begin{subproblems}
    \subproblem A set of polynomials $\{f_1, \ldots, f_k\}$ in $W$ is linearly
    independent if the equation
      \[c_1 f_1 + \cdots + c_k f_k = {\bf 0}\]
      has only the trivial solution $c_1 = \cdots c_k = 0$.

      In other words, if 
      \[c_1 f_1(x) + \cdots + c_k f_k(x) = 0\] for every real number $x$, then
      $c_1 = \cdots c_k = 0$.  %
      \subproblem A basis for $W$ is $\{f_1, f_2, f_3\}$ where
      \[f_1(x) = 1 + 0x + 0x^2, \qquad f_2(x) = 0 + 1x + 0x^2, 
      \qquad f_3(x) = 0 + 0x + 1x^2.\]
      
    First we show that $f_1, f_2$, and $f_3$ span $W$.  Let $f$ be any
    polynomial in $W$.  Then for some $a_0$, $a_1$, and $a_2$ we have
      \[f(x) = a_0 + a_1 x + a_2 x^2,\]
      so $f= a_0 f_1 + a_1 f_2 + a_2 f_3$ and therefore $f$ is in 
      $\textup{Span}\{f_1, f_2, f_3\}$.
      
      Also, $\{f_1, f_2, f_3\}$ is linearly independent: if 
      \[{\bf 0} = c_1 f_1 + c_2 f_2 + c_3 f_3,\]
      then for all $x$ we have
      \[0 = c_1 + c_2 x + c_3 x^2.\] This means that the polynomial
      $c_1 + c_2 x + c_3 x^2$ has infinitely many roots, which is only possible
      if all coefficients are $0$ (otherwise, it would have a maximum of $2$
      roots).  %
      \subproblem Since $W$ has a basis consisting of $3$ elements, the
      dimension of $W$ is $3$.
  \end{subproblems}
  \end{solution}
 \end{problems} \setcounter{problem}{0}

\newpage





\begin{center} \underline{2.9, The Rank Theorem and the Basis
    Theorem} \end{center}

\begin{problems}
\problem % Dan's notes
  True or false: There is a $3 \times 2$ matrix $A$ so that the column space of
  $A$ is a line and the null space of $A$ is a line.
  
\begin{solution}
  True: for example, $A = \mat{1 0; 1 0; 1 0}$.  This has one column with a
  pivot (so $\Col(A)$ is a line) and one column without a pivot (so $\Nul(A)$ is
  a line).
\end{solution}

\problem % Dan's notes
True or false: There is a $2 \times 3$ matrix $A$ so that the column space of
$A$ is a line and the null space of $A$ is a line. %

\begin{solution}
  False: If $\Col(A)$ and $\Nul(A)$ are both lines, then the sum of their
  dimensions is $2$, but this is impossible by the Rank Theorem since
\[\dim(\Col \ A) + \dim(\Nul \ A) = \ \textup{(number of columns of $A$)} = 3.\]
\end{solution}

\problem % Created for ILA
True or false: If $A$ is a
$10 \times 11$ matrix and its null space is a line, then every vector
in ${\bf R}^{10}$ is a linear combination of the columns of $A$.

\begin{solution}
  True.  First we recall that since $A$ has $10$ rows, the column
  space of $A$ is a subspace of ${\bf R}^{10}$. By the Rank Theorem
  and the fact that the null space is a line, we get
  \[\dim(\Col \ A) + \dim(\Nul \ A) = 11, \qquad
    \dim(\Col \ A) + 1 = 11, \qquad
    \dim(\Col \ A) = 10.\]
  Therefore, $\Col(A)$ is a $10$-dimensional subspace of ${\bf R}^{10}$,
  which is just another way of saying $\Col(A) = {\bf R}^{10}$.
\end{solution}

\problem % Created for ILA
Suppose $A$ is a $66 \times 77$ matrix.  Which of the following statements are true?  Select all that apply.
\begin{subproblems}
  \subproblem The homogeneous equation $Ax=0$ must have infinitely
  many solutions.
  \subproblem The column space of $A$ is a subspace ${\bf R}^{66}$.
  \subproblem $\dim(\Nul(A)) \geq 11$.
\end{subproblems}

\begin{solution}
  \begin{subproblems}
    \subproblem True: $A$ is $66 \times 77$, so $A$ has at most $66$
    pivots.  Since $A$ has $77$ columns, this means at least $11$
    columns in $A$ will not have pivots, therefore the equation
    $Ax=0$ has at least $11$ free variables.  This is just another
    way of rephrasing the consequences of the Rank Theorem, which
    expresses this succinctly in part (c).
    \subproblem True, since $A$ has $66$ rows.
    \subproblem True.  By the Rank Theorem,
    \[\dim(\Col \ A) + \dim(\Nul \ A) = 77,\]
    but since $\Col \ A$ is a subspace of ${\bf R}^{66}$ we know
    $\dim(\Col(A)) \leq 66$, therefore $\dim(\Nul \ A) \geq 11$.
  \end{subproblems}
\end{solution}

\problem
Is there a $3 \times 3$ matrix $A$ with the property 
that $\textup{Nul}(A) = \textup{Col}(A)$?  If so, write one.

\begin{solution}
No: if $\textup{Nul}(A) = \textup{Col}(A)$, then by
    the Rank Theorem,
    \[3 = \dim(\textup{Nul} \ A) + \dim(\textup{Col} \ A) = 2
      \dim(\textup{Nul} \ A).\]
    This would mean
    $\dim(\textup{Nul} \ A) = 3/2$.    
\end{solution}

\problem 
Is there a $2 \times 2$ matrix $A$ with the property 
that $\textup{Nul}(A) = \textup{Col}(A)$?  If so, write one.

\begin{solution}
Yes: unlike in the $3 \times 3$ case, the Rank Theorem
    does not rule out the possibility that
    $\textup{Nul}(A) = \textup{Col}(A)$.  However, it may take some
    time and trial and error to write such an $A$.  One example is
    $A = \mat{0 1; 0 0}$.  Both its range and null space are the span
    of $\vec{1 0}$.
\end{solution}


\problem % Dan's notes
Suppose that $A$ is a $6 \times 2$ matrix and that the column space of $A$ is
$2$-dimensional.  Is it possible for $\vec{1 0}$ and $\vec{1 1}$ to be solutions
to $Ax = b$ for some $b$ in ${\bf R}^6$?  Justify your answer.

\begin{solution}
  No, it is not possible.  Here $A$ is $6 \times 2$ with a $2$-dimensional
  column space, so both of its columns are pivot columns.  In other words, its
  columns are linearly independent.  Therefore, $Ax = 0$ has exactly one
  solution, so any consistent system $Ax = b$ must have exactly one solution.
\end{solution}

\problem % Dan's notes
Suppose that $A$ is a $5 \times 7$ matrix and the column space of $A$ is a line
in ${\bf R}^5$.  Describe the set of solutions to $Ax = 0$.

\begin{solution}
We are told that $\dim(\Col \ A) = 1$.  By the Rank Theorem:
\[\dim(\Col \ A) + \dim(\Nul \ A) = 7, \qquad
1 + \dim(\Nul \ A) = 7, \qquad \dim(\Nul \ A) = 6.\]
Also, the null space of $A$ lives in ${\bf R}^7$ since $A$ has $7$ columns.
Therefore, the set of solutions to $Ax = 0$ is $6$-dimensional subspace of ${\bf R}^7$.
\end{solution}

\problem % Dan's notes
Suppose that $A$ is a $5 \times 7$ matrix and the column space of $A$ is
${\bf R}^5$.  Describe the set of solutions to $Ax = 0$.

\begin{solution}
We are told that $\dim(\Col \ A) = 5$.  By the Rank Theorem:
\[\dim(\Col \ A) + \dim(\Nul \ A) = 7, \qquad
5 + \dim(\Nul \ A) = 7, \qquad \dim(\Nul \ A) = 2.\]
Also, the null space of $A$ lives in ${\bf R}^7$ since $A$ has $7$ columns.
Therefore, the set of solutions to $Ax = 0$ is a plane in ${\bf R}^7$.
\end{solution}

\problem % Dan's notes
Suppose that $A$ is a $5 \times 7$ matrix and the null space of $A$ is a plane.
Is the matrix equation $Ax = (1, 2, 3, 4, 5)$ consistent?  Justify your answer.

\begin{solution}
Yes.  The Rank Theorem gives us 
\[\dim(\Col \ A) + \dim(\Nul \ A) = 7, \qquad
  \dim(\Col \ A) + 2 = 7, \qquad \dim(\Col \ A) = 5.\] In other words, the
column space of $A$ is all of ${\bf R}^5$, therefore $Ax = b$ is consistent for
every $b$ in ${\bf R}^5$.
\end{solution}


\problem % Jankowski Spring 2018 exam 2
Fill in the blanks: If $A$ is a $5 \times 6$ matrix and whose columns span a
plane, then the null space of $A$ is a \underline{\qquad \qquad}-dimensional
subspace of %${\bf R}^{\underline{\qquad}}$.
$\R^{\fbox{\phantom{\tiny 2}}}$.

\begin{solution}
  $\Nul A$ is a subspace of ${\bf R}^6$, and
  $\dim(\Col \ A)+ \dim(\Nul \ A)=6$ so $2+\dim(\Nul A)=6$.  Thus
  $\Nul \ A$ is a 4-dimensional subspace of ${\bf R}^6$.
\end{solution}




\problem % Made for ILA
Write a matrix $A$ whose null space is $3$-dimensional and whose
column space is $4$-dimensional.

\begin{solution}
  By the Rank Theorem, the sum of the dimensions of the null and
  column spaces is equal to the number of columns of $A$, so $A$ must
  have $7$ columns.  Moreover, exactly $4$ columns must be pivot
  columns so that $\rank(A)=4$.

  One example is
\[A = \mat{1 0 0 0 0 0 0; 0 1 0 0 0 0 0; 0 0 1 0 0 0 0; 0 0 0 1 0 0 0}.\]
\end{solution}

\end{problems} \setcounter{problem}{0}


\newpage